% gif-server deployment guide.
% Source for deploy/deployment-guide.pdf. Rebuild with:
%   pdflatex -output-directory deploy deploy/deployment-guide.tex
% (run twice if section references ever get added).
\documentclass[11pt]{article}

% T1 + lmodern so literal $ < > _ { } | print correctly in \texttt (OT1 mangles
% them, e.g. $ -> pound sign).
\usepackage[T1]{fontenc}
\usepackage{lmodern}
\usepackage[a4paper,margin=1in]{geometry}
\usepackage{xcolor}
\usepackage{listings}
\usepackage{parskip}
\usepackage{framed}
\usepackage[hidelinks]{hyperref}
\usepackage{enumitem}

\definecolor{codebg}{rgb}{0.96,0.96,0.96}
\definecolor{codeframe}{rgb}{0.85,0.85,0.85}
\definecolor{warnbg}{rgb}{1.0,0.97,0.88}

% Block code style.
\lstdefinestyle{block}{
  backgroundcolor=\color{codebg},
  basicstyle=\ttfamily\small,
  breaklines=true,
  breakatwhitespace=false,
  frame=single,
  rulecolor=\color{codeframe},
  framesep=6pt,
  xleftmargin=4pt,
  xrightmargin=4pt,
  columns=fullflexible,
  keepspaces=true,
  showstringspaces=false,
  literate={~}{{\textasciitilde}}1
}
\lstset{style=block}

% Long monospace tokens (error strings, commands) can't hyphenate; let TeX
% stretch interword space rather than overrun the margin.
\sloppy

% Inline code shorthand. Uses \detokenize (not \lstinline) so it is safe inside
% macro arguments such as \colorbox{} and \item[...]; underscores, $ and # print
% literally without entering math mode.
\newcommand{\code}[1]{\texttt{\detokenize{#1}}}

% Simple "note" box: colored, breakable, accepts arbitrary content.
\colorlet{shadecolor}{warnbg}
\newenvironment{note}{\begin{shaded}\itshape\noindent\ignorespaces}{\end{shaded}}

\title{\textbf{gif-server --- Deployment Guide}}
\date{}

\begin{document}

{\Huge\bfseries gif-server --- Deployment Guide}\par
\medskip
{\small Self-hosted deployment to a Docker homeserver behind a containerized nginx
reverse proxy. Generated 2026-05-24.}

\hrule
\bigskip

\section*{1.\quad Architecture}
Three pieces, all on the homeserver:
\begin{itemize}[nosep]
  \item \textbf{db} --- Postgres 16 container (gif-server's compose stack).
  \item \textbf{app} --- the gif-server binary container. Migrations run automatically on
        startup. Bound to \code{127.0.0.1:<APP_HOST_PORT>} on the host (for debugging) and
        joined to the existing nginx Docker network with alias \code{gif-server}.
  \item \textbf{nginx} --- a pre-existing container (\code{lscr.io/linuxserver/nginx}) on the
        nginx network. The sole public ingress; terminates TLS and proxies to
        \code{http://gif-server:8847}.
\end{itemize}
The app is never published to the internet. The internet reaches it only through nginx.

\section*{2.\quad Prerequisites}
\begin{itemize}[nosep]
  \item Docker + Docker Compose on the homeserver (already present).
  \item The existing nginx container \emph{and the Docker network it is attached to}. Note the
        network's real name (see the box below) --- it is often \emph{not} literally \code{nginx}.
  \item DNS A/AAAA record for the public domain (e.g.\ \code{gifs.example.com}) pointing at the
        homeserver's public IP.
  \item Router port-forward of TCP \textbf{80} and \textbf{443} to the homeserver.
  \item A TLS certificate for the domain, mounted into the nginx container (handled separately).
  \item A few GB of free RAM/disk --- the first Rust build is slow and memory-hungry.
\end{itemize}

\begin{note}
\textbf{The nginx network name.} \code{docker-compose.yml} joins an \code{external} network. Compose
does a \emph{literal} lookup, so the network must exist by the exact name given. Having an nginx
\emph{container} is not the same as having a network \emph{named} \code{nginx}: a container started by
its own compose project is usually on a project-prefixed network like \code{nginx_default}. Find the
real name with \code{docker network ls} (or \code{docker inspect <nginx-container>}), then set it in the
\code{name:} field of the network block --- see Section 7.
\end{note}

\section*{3.\quad Repository files involved}
\begin{description}[style=nextline,nosep,leftmargin=0pt]
  \item[\code{docker-compose.yml}] db + app stack. Config comes from \code{.env} via
        \code{${VAR}} interpolation. App joins the external nginx network.
  \item[\code{.env.example}] Template for configuration. The real \code{.env} is created on the
        server and is git-ignored / not rsynced.
  \item[\code{deploy/nginx.conf.example}] Reverse-proxy server block to copy into the nginx
        container's site-confs.
  \item[\code{deploy/deployment-guide.tex}] Source for this PDF. Edit it and recompile rather than
        editing the PDF directly.
  \item[\code{Dockerfile}] Multi-stage Rust build (cargo-chef) producing a slim runtime image.
\end{description}

\section*{4.\quad Step 1 --- Copy the code to the server}
From the development machine, in the project directory:
\begin{lstlisting}
rsync -avz --delete \
  --exclude target/ --exclude .git/ --exclude /data/ --exclude .env \
  ./ user@homeserver:/srv/gif-server/
\end{lstlisting}
Excludes the build cache, git history, server-side GIF data, and the env file. Run once with
\code{-n} (dry run) first to preview, then without.

\begin{note}
\textbf{Anchor the data-dir exclude with a leading slash} (\code{/data/}, not \code{data/}). An
unanchored pattern matches a directory of that name \emph{anywhere} in the tree --- which is how an
earlier \code{--exclude storage/} silently dropped the \code{src/storage/} \emph{module} and broke the
build (\code{error[E0583]: file not found for module storage}). The runtime data dir is named
\code{data} precisely so it cannot collide with any source directory.
\end{note}

First time: ensure the target exists and is writable by your SSH user:
\begin{lstlisting}
sudo mkdir -p /srv/gif-server
sudo chown "$USER:$USER" /srv/gif-server
\end{lstlisting}

\section*{5.\quad Step 2 --- Create the production \texttt{.env} on the server}
At \code{/srv/gif-server/.env} (it is not rsynced, so create it fresh). Generate secrets with
\code{openssl rand}:
\begin{lstlisting}
POSTGRES_PASSWORD=<openssl rand -hex 24>
SESSION_SECRET=<openssl rand -hex 32>
BASE_URL=https://gifs.example.com
MATRIX_SERVER_NAME=matrix.example.com
MATRIX_FEDERATION_URL=https://matrix.example.com
MATRIX_ADMIN_MXIDS=@admin:matrix.example.com
CORS_ALLOWED_ORIGINS=https://<your-frontend-origin>
# Optional: override host-published ports only if they clash with other
# host services (the app/nginx use the internal network, not these).
DB_HOST_PORT=5433
APP_HOST_PORT=8848
\end{lstlisting}
\begin{itemize}[nosep]
  \item \code{DATABASE_URL} is \textbf{not} needed here --- compose builds it from
        \code{POSTGRES_PASSWORD} with host \code{db}.
  \item \code{SESSION_SECRET} must be at least 32 characters.
  \item \code{BASE_URL} must be the public \code{https://} URL --- it is embedded in rendition URLs.
  \item \code{MATRIX_FEDERATION_URL} is where the homeserver's federation API is reachable. If the
        server delegates federation to port 443 via \code{.well-known/matrix/server}, this is just
        \code{https://<server-name>} (no \code{:8448}).
  \item \code{DB_HOST_PORT} / \code{APP_HOST_PORT} set the host ports the db and app containers
        publish on \code{127.0.0.1} (defaults \code{5432} / \code{8847}). Nothing in the running
        system depends on them --- the app reaches Postgres at \code{db:5432} and nginx reaches the
        app at \code{gif-server:8847}, both over the internal network --- so override them only to
        dodge a host-side port clash. If you change \code{APP_HOST_PORT}, adjust the host-side
        \code{curl .../health} checks accordingly.
  \item \code{STORAGE_MAX_BYTES} / \code{PER_USER_STORAGE_BYTES} fall back to defaults (10GB / 1GB)
        if omitted.
\end{itemize}

\begin{note}
Create \code{.env} \textbf{before} the first \code{docker compose up}. Required variables use
\code{${VAR:?}} markers, so compose fails loudly (by design) if any are missing.
\end{note}

\section*{6.\quad Step 3 --- Build and start the stack}
\begin{lstlisting}
cd /srv/gif-server
docker compose up -d --build
docker compose logs -f app          # watch migrations apply + startup
curl -fsS http://127.0.0.1:8847/health   # expect HTTP 200 from the host
\end{lstlisting}
A clean boot logs the migrations (\code{001_init.sql}, \code{002_hidden.sql}) then
``listening on 0.0.0.0:8847''.

\section*{7.\quad Step 4 --- Wire up the nginx network and site}
\textbf{Network.} If \code{docker compose up} fails with
\code{network nginx declared as external, but could not be found}, point the compose network block
at the real network name reported by \code{docker network ls}:
\begin{lstlisting}
# docker-compose.yml
networks:
  default:
  nginx:
    external: true
    name: nginx_default     # <-- the actual network the nginx container is on
\end{lstlisting}
Alternatively, create a network literally named \code{nginx} and attach the container:
\begin{lstlisting}
docker network create nginx
docker network connect nginx <nginx-container>
\end{lstlisting}

\textbf{Site.} Copy \code{deploy/nginx.conf.example} into the nginx container's site-confs directory
on the host:
\begin{lstlisting}
cp /srv/gif-server/deploy/nginx.conf.example \
  /docker/apps/nginx/nginx/site-confs/gif-server.conf
\end{lstlisting}
Edit it: set the real \code{server_name} and the container-internal TLS cert paths. Key points
already baked into the template:
\begin{itemize}[nosep]
  \item Proxies to \code{http://gif-server:8847} using Docker's embedded DNS resolver
        (\code{127.0.0.11}) and a variable upstream, so an app restart doesn't strand nginx on a
        stale IP.
  \item \code{client_max_body_size 50m;} --- matches the app's 50 MB upload limit (nginx defaults to
        1 MB).
  \item Overwrites \code{X-Forwarded-For} / \code{X-Real-IP} with \code{$remote_addr}.
\end{itemize}
Then validate and reload nginx:
\begin{lstlisting}
docker exec <nginx-container> nginx -t
docker exec <nginx-container> nginx -s reload
curl -fsS https://gifs.example.com/health   # expect HTTP 200 over TLS
\end{lstlisting}

\section*{8.\quad Step 5 --- Verify}
\begin{itemize}[nosep]
  \item \textbf{App resolvable from nginx:} \code{docker exec <nginx-container> getent hosts gif-server}
        returns an IP.
  \item \textbf{Real client IPs:} nginx access logs show public client IPs, not \code{172.x} gateway
        addresses --- confirms the per-IP rate limiter keys on real clients.
  \item \textbf{Auth path:} \code{POST /auth/matrix} with a real Matrix OpenID token returns a JWT;
        bad tokens return 401 (not 500).
\end{itemize}

\begin{note}
\textbf{Security --- X-Forwarded-For.} The auth rate limiter trusts the leftmost
\code{X-Forwarded-For} value. nginx must \textbf{overwrite} it with \code{$remote_addr} (the template
does), \textbf{not} append via \code{$proxy_add_x_forwarded_for} --- appending would let a client
prepend a spoofed IP and bypass the per-IP login limiter. The app must remain reachable only through
nginx (it is: localhost bind + internal network, never publicly published). A non-spoofable global
auth cap ($\sim$5 req/s) bounds total login load as a second layer.
\end{note}

\section*{9.\quad Local development (unchanged)}
Local dev runs the app on the host against the dockerized db:
\begin{lstlisting}
./scripts/dev.sh          # starts db via compose, waits healthy, runs cargo run
# or, if db is already up:
cargo run
curl -i http://localhost:8847/health
\end{lstlisting}
\begin{itemize}[nosep]
  \item The app reads the local \code{.env} (localhost DB URL, localhost base URL).
  \item \code{docker compose up -d db} works locally despite the external nginx network (only the
        app service references it).
  \item Running the \textbf{full} stack locally (\code{docker compose up}) requires the network to
        exist: \code{docker network create nginx} first.
\end{itemize}

\section*{10.\quad Routine operations}
\begin{description}[style=nextline,nosep,leftmargin=0pt]
  \item[Redeploy after code change] rsync (Step 1), then \code{docker compose up -d --build}
  \item[View app logs] \code{docker compose logs -f app}
  \item[Restart app only] \code{docker compose restart app}
  \item[Stop everything] \code{docker compose down} (data persists in the postgres\_data volume)
\end{description}

\begin{note}
GIF files live in the host bind mount \code{/srv/gif-server/data} and are root-owned (the container
runs as root). Back up that directory plus the Postgres volume.
\end{note}

\section*{11.\quad Troubleshooting}
\begin{description}[style=nextline,leftmargin=0pt]
  \item[\code{error[E0583]: file not found for module storage} (plus a cascade of \code{E0282})]
        The \code{src/storage/} module never reached the server. Cause: an unanchored rsync exclude
        (\code{storage/}) matching the source dir. Fix: use \code{--exclude /data/} (the data dir is
        named \code{data}, so it no longer collides), and confirm \code{src/storage/} is present on
        the server before building.
  \item[\code{Bind for 0.0.0.0:5432 failed: port is already allocated}]
        Another Postgres on the host already owns the published port. Set \code{DB_HOST_PORT} in the
        server's \code{.env} to a free port (e.g.\ 5433). The app is unaffected --- it uses the
        internal \code{db:5432}, not the host port.
  \item[\code{network nginx declared as external, but could not be found}]
        No Docker network with that exact name exists. Set the \code{name:} field of the compose
        network block to the real network (Section 7), or create/attach a network literally named
        \code{nginx}.
\end{description}

\end{document}
